\documentclass[runningheads]{llncs}

\usepackage{amssymb,amsmath,mathrsfs}
\usepackage[small,basic]{complexity}
\usepackage{xspace}
\usepackage{mathtools}
\usepackage{thmtools}

\usepackage{thm-restate}
\usepackage{setspace}



\usepackage[dvipsnames,table]{xcolor}
\definecolor{ruby}{RGB}{116,0,1}
\definecolor{emerald}{RGB}{26,121,42}
\definecolor{topaz}{RGB}{236,185,57}
\definecolor{sapphire}{RGB}{14,26,164}
\usepackage{tikz}


\newcommand{\red}[1]{\textcolor{ruby}{#1}}
\newcommand{\blue}[1]{\textcolor{sapphire}{#1}}


\usepackage{bm}
%
\usepackage{enumerate}


\usepackage{hyperref}
\hypersetup{
	final,
	colorlinks,
	 linkcolor={sapphire},
	 citecolor={emerald},
	 urlcolor={ruby}
}



\usepackage[capitalise,nameinlink]{cleveref}
\newcommand{\appref}[1]{\hyperref[#1]{Appendix~\ref*{#1}}}


%


\usepackage{listings} 
\lstdefinelanguage{Sage}[]{Python}
{morekeywords={True,False,sage,singular},
sensitive=true}
\lstset{frame=single,
          breaklines=true,
          postbreak=\mbox{\textcolor{red}{$\hookrightarrow$}\space},
          showtabs=False,
          showspaces=False,
          showstringspaces=False,
          commentstyle={\ttfamily\color{dredcolor}},
          keywordstyle={\ttfamily\color{dbluecolor}\bfseries},
          stringstyle ={\ttfamily\color{dgraycolor}\bfseries},
          language = Sage,
	  basicstyle={\small \ttfamily},
	  aboveskip=10pt,
	  belowskip=10pt,
          }
\definecolor{dblackcolor}{rgb}{0.0,0.0,0.0}
\definecolor{dbluecolor}{rgb}{.01,.02,0.7}
\definecolor{dredcolor}{rgb}{0.8,0,0}
\definecolor{dgraycolor}{rgb}{0.30,0.3,0.30}
\newcommand{\dblue}{\color{dbluecolor}\bf}
\newcommand{\dred}{\color{dredcolor}\bf}
\newcommand{\dblack}{\color{dblackcolor}\bf}


\newcommand{\Sage}{{\color{dbluecolor}\sf Sage}\xspace}


\newcommand{\keywords}[1]{\par\addvspace\baselineskip
\noindent\keywordname\enspace\ignorespaces#1}



% Theorem-like environments.
\spnewtheorem{nclam}{Claim}{\bfseries}{\itshape}
\crefname{nclam}{Claim}{Claims}
% claims in separate counter

\let\remark\relax % undefine the environment
\let\theremark\relax
\spnewtheorem{remark}{Remark}{\bfseries}{\rmfamily} % puts remarks into separate counter

%%%%%%

\newcommand{\seq}[3][{}]{\langle #2 \rangle_{#3}^{#1}}

% Complexity classes
\newclass{\EXPTIME}{EXPTIME}
\newclass{\PTIME}{PTIME}

% Shorthand commands
\renewcommand{\R}{\mathbb{R}}
\renewcommand{\C}{\mathbb{C}}
\newcommand{\Q}{\mathbb{Q}}
\newcommand{\QQ}{\mathbb{Q}}
\newcommand{\Z}{\mathbb{Z}}
\newcommand{\N}{\mathbb{N}}
\newcommand{\B}{\mathcal{B}}
\newcommand{\Oc}{\mathcal{O}}
\renewcommand{\H}{\mathbb{H}}
\renewcommand{\K}{\mathbb{K}}

\DeclareMathOperator{\GL}{GL}
\DeclareMathOperator{\PGL}{PGL}
\DeclareMathOperator{\PSL}{PSL}
\let\SL\relax \DeclareMathOperator{\SL}{SL}
\DeclareMathOperator{\SU}{SU}
\DeclareMathOperator{\SZ}{SZ}
\DeclareMathOperator{\SA}{SA}
\DeclareMathOperator{\SM}{SM}
\DeclareMathOperator{\SO}{SO}
\DeclareMathOperator{\T}{T}
\DeclareMathOperator{\Id}{Id}
\DeclareMathOperator{\Lat}{Lat}
\DeclareMathOperator{\uni}{UT}
\DeclareMathOperator{\ut}{U}
\DeclareMathOperator{\fg}{F}
\DeclareMathOperator{\sg}{S}

\renewcommand{\Re}{\operatorname{Re}}
\renewcommand{\Im}{\operatorname{Im}}
\newcommand{\veca}{\bm{a}}
\newcommand{\vecb}{\bm{b}}
\newcommand{\vecz}{\bm{0}}
%\newcommand{\vect}{\textrm{vec}}

%\usepackage{algpseudocode}
\usepackage[plain]{algorithm}
\usepackage{algorithmicx,algpseudocode}


 \renewcommand{\algorithmicrequire}{\textbf{Input:}}
 \renewcommand{\algorithmicensure}{\textbf{Assert:}}
 \renewcommand{\algorithmiccomment}[1]{\hfill\textcolor{gray}{$\triangleright$ #1}}


\DeclareMathOperator{\tr}{tr}

\newcommand*\wc{{ }\cdot{ }}

\renewcommand{\S}{\mathcal{S}}

\newcolumntype{L}[1]{>{\centering\arraybackslash}m{#1}}


% The number `e'
\providecommand*{\eu}%
{\ensuremath{\mathrm{e}}}
% The imaginary unit
\providecommand*{\iu}%
{\ensuremath{\mathrm{i}}}


% short minus sign

\DeclareMathSymbol{\shortminus}{\mathbin}{AMSa}{"39}
\newcommand{\medminus}{\scalebox{0.6}[0.7]{\(-\)}}
\newcommand{\minus}{\mathchoice{-}{-}{\medminus}{\shortminus}}


\usepackage[left=25mm, right=25mm]{geometry}

\usepackage[T1]{fontenc}
\usepackage[utf8]{inputenc} \DeclareUnicodeCharacter{2212}{-}


\begin{document}
\title{SageMath Code for: On Word Representations and Embeddings in Complex Matrices
}
\titlerunning{On Word Representations and Embeddings in Complex Matrices}

\author{George Kenison}
\institute{}

\maketitle

\section{Pseudocode}

In the appendix of our paper, we prove a claim by employing an exhaustive search of matrices in \(\PSL(2,\Oc_d)\) with a specified norm property. The pseudocode below details the procedure.

\begin{algorithm}
\setstretch{1.35}
\begin{algorithmic}
\Require {\(d\in\{1,2,3,7,11\}\)}
\State \texttt{matrices} \(\leftarrow\) empty list 
\State \(S\leftarrow \bigl\{x\in \Oc_d : N(x) < \tfrac{1}{1-\kappa(d)}\bigr\}\); \Comment{set of matrix entries (see full paper)}
\For{ordered tuples \((\alpha,\beta,\gamma, \delta)\), with \(\alpha,\beta,\gamma,\delta\in S\), \(\gamma\neq 0\), and \(\alpha \delta - \beta \gamma = 1\)} 
\State \(\bigl( M, \|M\|\bigr)  \leftarrow \left( \begin{psmallmatrix} \alpha & \beta \\ \gamma & \delta \end{psmallmatrix} , \max\{|\alpha|^2, |\beta|^2, |\gamma|^2, |\delta|^2\}\right)\); \Comment{matrix and its norm}
\State \(\theta \leftarrow \textsc{Quotient}(-\delta, \gamma)\); \Comment{output quotient \(-\theta\in\Oc_d\)} %
\State \((\alpha,\beta) \leftarrow (\theta \cdot \alpha + \beta, -\alpha)\); 
\State \((\gamma, \delta) \leftarrow (\theta \cdot \gamma + \delta, -\gamma)\);
\State \(\bigl( N, \|N\| \bigr)  \leftarrow \left(\begin{psmallmatrix} \alpha & \beta \\ \gamma & \delta \end{psmallmatrix}, \max\{|\alpha|^2, |\beta|^2, |\gamma|^2, |\delta|^2\} \right)  \); \Comment{update matrix and its norm}
\If{\(\|M\| < \|N\|\)}
\State add \((M, N)\) to \texttt{matrices} %
\EndIf
\EndFor
\If{\texttt{matrices} is an empty list}
\State \Return {\lq\lq Claim holds for matrices in \(\PSL(2,\Oc_d)\)\rq\rq}
\Else
\State \Return {\lq\lq Claim does not hold for the pairs} \texttt{matrices}\rq\rq
\EndIf
\caption{Search and verification of claim in \(\PSL(2,\Oc_d)\).\label{alg:cap}}
\end{algorithmic}
\end{algorithm}

In \cref{alg:cap}, the function \(\textsc{Quotient}\colon \Oc_d \times \Oc_d \to \Oc_d\) takes as input an ordered pair of algebraic integers and returns their quotient (in the ring \(\Oc_d\)).
In other words, for the ordered pair \((a,b)\in\Oc_d\times \Oc_d\), \textsc{Quotient} returns the nearest integer (in \(\Oc_d\)) to \(a/b\).

\section{SageMath code for \texorpdfstring{\(\PSL(2,\Z[\iu])\)}{PSL(2,Z[i])}}
Here we comment on the implementation for the search for matrices in \(\PSL(2,\Z[\iu])\). \textit{Mutatis mutandis,} the procedures for searching the other Euclidean Bianchi groups are similar.


\medskip
We define the Gaussian field \(\QQ[\iu]\) and fix the embedding that maps the generator \(\omega\), a solution of \(\omega^2+1=0\), to \(\omega \mapsto \iu\).  We also specify the Gaussian integers \(\Z[\iu]\).
%
\begin{lstlisting}
sage: # CC.0 is the imaginary unit 'i'
      K.<w> = QuadraticField(-1, embedding=CC.0)
      OK = K.ring_of_integers()
\end{lstlisting}

SageMath does not automatically recognise that the integer rings \(\Oc_d\) with \(d=1,2,3,7,11\) are Euclidean domains.
Thus for integers \(a,b\in\Oc_d\), the SageMath command \verb|a.quo_rem(b)[0]| outputs the result of the division \(a/b\in\Q(\sqrt{-d})\) rather than the nearest integer in \(\Oc_d\) (the output of \textsc{Quotient}).
We circumvent this obstacle by implementing a straightforward procedure for finding the nearest algebraic integer in \(\Oc_d\).

As an aside, simply rounding the coordinates of an algebraic number does not work for finding the nearest algebraic integer in an oblique lattice (e.g., the nearest integer in \(\Oc_7\)).
\begin{lstlisting}
sage: # Function that returns the nearest algebraic integer in OK
      def nearest_integer(field_quo):
        K = field_quo.parent()
    
        # {1, omega} is a basis for the integer lattice
        basis = OK.basis()
    
        # Embed the integer lattice and field_quo into RR^2. field_quo maps to field_target
        def embed(element):
            v = CC(element) 
            return vector(RR, [v.real(), v.imag()])
        field_target = embed(field_quo)
    
        # Basis matrix B for the integer lattice
        B = matrix(RR, [embed(b) for b in basis]).transpose()
    
        # Exact coordinates of the field_target in the lattice basis
        coords = B.solve_right(field_target)
    
        # Variables that record the current nearest integer
        best_norm = infinity
        best_integer = None

        # Search local neighbours for the nearest integer to field_target
        import itertools
        options = [[floor(v)-1, floor(v), floor(v)+1] for v in coords]
        for p, q in itertools.product(*options):
            candidate = OK([p, q])
            # Norm of diff between candidate and field_target
            diff = embed(candidate) - field_target
            norm_diff = diff.dot_product(diff)

            # Update current nearest integer 
            if norm_diff < best_norm:
                best_norm = norm_diff
                best_integer = candidate
            
      return best_integer
\end{lstlisting}

Now we generate all \(2\times 2\) matrices with entries in \(\Z[\iu]\) and determinant \(1\).  We also produce the list of updated matrices, where we employ the \verb|nearest_integer()| function defined above.

Recall that a matrix \(M\) and its update \(N\) in \(\PSL(2,\Oc_d)\) are defined as follows:
\begin{equation*}M = \begin{psmallmatrix} \alpha & \beta \\ \gamma & \delta \end{psmallmatrix} \quad \text{and} \quad
%\label{eq:update}
    N =
    \begin{psmallmatrix}
    \theta \alpha + \beta & -\alpha \\ 
    \theta \gamma + \delta & - \gamma
    \end{psmallmatrix}.
\end{equation*}
\begin{lstlisting}
sage: # Matrix entries in ZZ[i] (for other OK, see full paper or implementation)
      S = [K(0), K(1), w, w*w, w*w*w]
      T = [K(1), w, w*w, w*w*w]

      # Generate lists of matrices and their updates
      mats = []
      matsupdate = []

      for a in S:
        for b in S:
            for c in T: 
                for d in S:
                    M = matrix(K, [[a, b], [c, d]])
                        if M.det() == 1:
                            # nearest integer to quotient
                            theta = nearest_integer(K(K(-d).quo_rem(K(c))[0]))
                            # update entry a
                            alpha = K(theta)*K(a) + K(b)
                            # update entry c
                            gamma = K(theta)*K(c) + K(d)
                            N = matrix(K, [[alpha, -a], [gamma, -c]])                        
                            mats.append(M)                 
                            matsupdate.append(N)
\end{lstlisting}

We compute the norms for the matrices and their updates.
\begin{lstlisting}
sage: # Return the max of the normed elements in a matrix
      def max_abs_squared_norm(M):
        return max((x*x.galois_conjugate() for x in M.list()))

      # Generate lists for the norms of the matrices and their updates
      normmats = []
      normmatsupdate = []
      
      for m in mats:
        normmats.append(max_abs_squared_norm(m))

      for m in matsupdate:
        normmatsupdate.append(max_abs_squared_norm(m))
\end{lstlisting}

Output the list of matrices, updates, and their respective norms.
\begin{lstlisting}
sage: result = list(zip(mats, matsupdate, normmats, normmatsupdate))

      # print the matrices, their updates, and the respective norms
      for i, (M, N, x, y) in enumerate(result, 1):
        print(f"#{i:3d} matrix {list(map(list, M.rows()))} update {list(map(list, N.rows()))} | norms=({x}, {y})")
\end{lstlisting}


%\begin{minipage}[c]{0.95\textwidth}
\begin{lstlisting}[caption={Output candidate matrices, update matrices, and their respective norms. \label{list:here}}]
  1. matrix [[0, 1], [-1, 0]] update [[1, 0], [0, 1]] | norms=(1, 1)
  2. matrix [[0, 1], [-1, 1]] update [[1, 0], [0, 1]] | norms=(1, 1)
  3. matrix [[0, 1], [-1, w]] update [[1, 0], [0, 1]] | norms=(1, 1)
  4. matrix [[0, 1], [-1, -1]] update [[1, 0], [0, 1]] | norms=(1, 1)
  5. matrix [[0, 1], [-1, -w]] update [[1, 0], [0, 1]] | norms=(1, 1)
  6. matrix [[0, w], [w, 0]] update [[w, 0], [0, -w]] | norms=(1, 1)
  7. matrix [[0, w], [w, 1]] update [[w, 0], [0, -w]] | norms=(1, 1)
  8. matrix [[0, w], [w, w]] update [[w, 0], [0, -w]] | norms=(1, 1)
  9. matrix [[0, w], [w, -1]] update [[w, 0], [0, -w]] | norms=(1, 1)
 10. matrix [[0, w], [w, -w]] update [[w, 0], [0, -w]] | norms=(1, 1)
 11. matrix [[0, -1], [1, 0]] update [[-1, 0], [0, -1]] | norms=(1, 1)
 12. matrix [[0, -1], [1, 1]] update [[-1, 0], [0, -1]] | norms=(1, 1)
 13. matrix [[0, -1], [1, w]] update [[-1, 0], [0, -1]] | norms=(1, 1)
 14. matrix [[0, -1], [1, -1]] update [[-1, 0], [0, -1]] | norms=(1, 1)
 15. matrix [[0, -1], [1, -w]] update [[-1, 0], [0, -1]] | norms=(1, 1)
 16. matrix [[0, -w], [-w, 0]] update [[-w, 0], [0, w]] | norms=(1, 1)
 17. matrix [[0, -w], [-w, 1]] update [[-w, 0], [0, w]] | norms=(1, 1)
 18. matrix [[0, -w], [-w, w]] update [[-w, 0], [0, w]] | norms=(1, 1)
 19. matrix [[0, -w], [-w, -1]] update [[-w, 0], [0, w]] | norms=(1, 1)
 20. matrix [[0, -w], [-w, -w]] update [[-w, 0], [0, w]] | norms=(1, 1)
 21. matrix [[1, 0], [1, 1]] update [[-1, -1], [0, -1]] | norms=(1, 1)
 22. matrix [[1, 0], [w, 1]] update [[w, -1], [0, -w]] | norms=(1, 1)
 23. matrix [[1, 0], [-1, 1]] update [[1, -1], [0, 1]] | norms=(1, 1)
 24. matrix [[1, 0], [-w, 1]] update [[-w, -1], [0, w]] | norms=(1, 1)
 25. matrix [[1, 1], [-1, 0]] update [[1, -1], [0, 1]] | norms=(1, 1)
 26. matrix [[1, w], [w, 0]] update [[w, -1], [0, -w]] | norms=(1, 1)
 27. matrix [[1, -1], [1, 0]] update [[-1, -1], [0, -1]] | norms=(1, 1)
 28. matrix [[1, -w], [-w, 0]] update [[-w, -1], [0, w]] | norms=(1, 1)
 29. matrix [[w, 0], [1, -w]] update [[-1, -w], [0, -1]] | norms=(1, 1)
 30. matrix [[w, 0], [w, -w]] update [[w, -w], [0, -w]] | norms=(1, 1)
 31. matrix [[w, 0], [-1, -w]] update [[1, -w], [0, 1]] | norms=(1, 1)
 32. matrix [[w, 0], [-w, -w]] update [[-w, -w], [0, w]] | norms=(1, 1)
 33. matrix [[w, 1], [-1, 0]] update [[1, -w], [0, 1]] | norms=(1, 1)
 34. matrix [[w, w], [w, 0]] update [[w, -w], [0, -w]] | norms=(1, 1)
 35. matrix [[w, -1], [1, 0]] update [[-1, -w], [0, -1]] | norms=(1, 1)
 36. matrix [[w, -w], [-w, 0]] update [[-w, -w], [0, w]] | norms=(1, 1)
 37. matrix [[-1, 0], [1, -1]] update [[-1, 1], [0, -1]] | norms=(1, 1)
 38. matrix [[-1, 0], [w, -1]] update [[w, 1], [0, -w]] | norms=(1, 1)
 39. matrix [[-1, 0], [-1, -1]] update [[1, 1], [0, 1]] | norms=(1, 1)
 40. matrix [[-1, 0], [-w, -1]] update [[-w, 1], [0, w]] | norms=(1, 1)
 41. matrix [[-1, 1], [-1, 0]] update [[1, 1], [0, 1]] | norms=(1, 1)
 42. matrix [[-1, w], [w, 0]] update [[w, 1], [0, -w]] | norms=(1, 1)
 43. matrix [[-1, -1], [1, 0]] update [[-1, 1], [0, -1]] | norms=(1, 1)
 44. matrix [[-1, -w], [-w, 0]] update [[-w, 1], [0, w]] | norms=(1, 1)
 45. matrix [[-w, 0], [1, w]] update [[-1, w], [0, -1]] | norms=(1, 1)
 46. matrix [[-w, 0], [w, w]] update [[w, w], [0, -w]] | norms=(1, 1)
 47. matrix [[-w, 0], [-1, w]] update [[1, w], [0, 1]] | norms=(1, 1)
 48. matrix [[-w, 0], [-w, w]] update [[-w, w], [0, w]] | norms=(1, 1)
 49. matrix [[-w, 1], [-1, 0]] update [[1, w], [0, 1]] | norms=(1, 1)
 50. matrix [[-w, w], [w, 0]] update [[w, w], [0, -w]] | norms=(1, 1)
 51. matrix [[-w, -1], [1, 0]] update [[-1, w], [0, -1]] | norms=(1, 1)
 52. matrix [[-w, -w], [-w, 0]] update [[-w, w], [0, w]] | norms=(1, 1)

\end{lstlisting}
%\end{minipage}

We automate verifying the claim in the following code.  The claim is checked for all candidate matrix pairs arising from \texttt{mats} and \texttt{matsupdate}.  For \(\PSL(2,\Z[\iu])\) the matrix pairs are given in Listing~\ref{list:here}.
\begin{lstlisting}
sage: # test the claim
      claimmats = []
      for i, (M, N, x, y) in enumerate(result, 1):
        if x < y:
            print(f"#{i:3d} matrix {list(map(list, M.rows()))} update {list(map(list, N.rows()))} | norms=({x}, {y})")
            claimmats.append(M)  
        
      if not claimmats:
        print("claim holds")
    
      if claimmats:
        print ("claim does not hold")
\end{lstlisting}

\end{document}